\vspace*{2cm}
\begin{center}
\textbf{\large Abstract}
\end{center}

\noindent
Continuous monitoring of radiological pollution relies on environmental
surveillance beacons that record the rate of gamma photons reaching a
scintillation detector. The recorded signal is heavily contaminated by
sensor noise and slow baseline drift, and the events of interest ---
spikes, plateaus, ramp-ups, ramp-downs --- are sparse and exhibit a
wide range of shapes, amplitudes and durations. The goal of this
\emph{Pôle Projet} at CentraleSup\'elec was to design, train, and
evaluate deep neural networks that can classify these events on
streaming gamma-dose recordings.

The project ran in two iterations. The \emph{first} (v1) calibrated a
white-Gaussian-plus-Ornstein--Uhlenbeck background model from four
months of real 2015 recordings and injected six hand-designed anomaly
templates densely into a $10^{7}$-sample synthetic signal; three
PyTorch architectures (a 1D-CNN, a 1D residual network, and a hybrid
CNN+Transformer) were trained on this dataset. On the held-out
synthetic test segment, all three models cleared $90\%$ event-level
macro F1, with the residual network reaching a flawless $100\%$. On
the \emph{real} 2015 data, however, all three collapsed in different
ways --- continuous false-positive plateaus, oscillating predictions,
single-class fixation. The synthetic--real distribution shift turned
out to be much larger than the noise model had suggested.

The \emph{second} iteration (v2) re-analysed the background process,
discovered that the high-frequency residual is well described by an
AR(1) Gaussian rather than IID white noise ($\varphi \approx 0.86$,
$\sigma_\varepsilon \approx 1.12$), hand-labelled the seventeen
plausible events visible across the four months, and re-tuned the
generator to (i) reproduce the AR(1) noise texture, (ii) match the
empirical anomaly amplitude and duration ranges, (iii) lower the
injection density from $\sim 7$ to $\sim 1.5$ events per $10\text{k}$
samples, and (iv) prune the template family to the three shapes
(\texttt{bell}, \texttt{fast\_ascend}, \texttt{fast\_descend}) that
actually occur in the labels. Retrained on this corpus, the same three
architectures now \emph{commit} to the normal class on quiet stretches
of the real signal and fire only where a visible event is present.

This report documents both iterations: the v1 retrospective, the
re-analysis that exposed its blind spots, the design of v2, and the
v2 results on synthetic data, robustness sweeps, and real recordings.

\vfill
\begin{center}
\small
\textit{Keywords:} time series classification, 1D convolutional neural
networks, residual networks, self-attention, transformer encoder,
radiological monitoring, synthetic data generation, AR(1) noise,
distribution shift.
\end{center}
